\documentclass[11pt,a4paper]{article}

\usepackage[T1]{fontenc}
\usepackage[utf8]{inputenc}
\usepackage{lmodern}
\usepackage[margin=2.5cm]{geometry}
\usepackage{microtype}
\usepackage{booktabs}
\usepackage{amsmath,amssymb}
\usepackage{enumitem}
\usepackage[hidelinks]{hyperref}

\title{PODS 2025 --- Conference Report}
\author{Group seminar, Algorithms course, UL FRI \\ \small{(authors to be filled in)}}
\date{\today}

\begin{document}
\maketitle

\section{The conference at a glance}

The 44th ACM SIGMOD-SIGACT-SIGAI Symposium on Principles of Database Systems
(PODS 2025) was held jointly with SIGMOD 2025 in Berlin, Germany, in June 2025.
The proceedings appear as Volume~3, Issue~2 of the \emph{Proceedings of the ACM
on Management of Data} (PACMMOD), published in May 2025
(\url{https://dl.acm.org/toc/pacmmod/2025/3/N2}). The volume opens with an
editorial by Carmeli, Geerts, Kara and Kimelfeld (Article~87) and contains
30~research papers as Articles~88--117.

This report covers the December~2024 submission cycle of PODS~2025, whose
papers populate the issue above. Out of $80$ submitted papers, $30$ were
accepted, an acceptance rate of $37.5\%$. (For context, across both
submission cycles the conference accepted $45$ of $127$ papers, $35.4\%$, but
the second cycle is published in a separate issue.)

\subsection*{Topics}

PODS is the flagship venue for theoretical work on database systems. The 30
papers in this issue cluster around the following themes:

\begin{itemize}[leftmargin=*,itemsep=0pt]
  \item query languages and their evaluation: conjunctive queries, regular
        path queries, Datalog, and disjunctive variants;
  \item information-theoretic and worst-case-optimal join algorithms,
        including new uses of fast matrix multiplication;
  \item differential privacy, streaming and sketching algorithms over
        relational data, including robust variants over noisy inputs;
  \item randomised data structures and their adversarial robustness, in
        particular locality-sensitive hashing and sketches;
  \item graph algorithms, dynamic data structures and complexity theory;
  \item explainability, provenance and responsibility measures for query
        answers;
  \item data integration, schema matching and emerging applications such as
        quantum optimisation for database problems.
\end{itemize}

\section{Selected paper summaries}

We selected three papers spanning very different parts of this landscape:
a deep theoretical advance in worst-case-optimal join evaluation that
folds fast matrix multiplication into the polymatroid framework; a sharp
attack on locality-sensitive hashing that quantifies how fragile its
textbook guarantee becomes against an adaptive adversary; and a streaming
algorithm for robust distinct-element counting and $\ell_0$-sampling in
general metric spaces, with a matching lower bound.

\subsection*{Abo Khamis, Hu, Suciu.\\
\emph{Fast Matrix Multiplication meets the Submodular Width.}\\
PACMMOD Vol.~3, No.~2 (PODS~'25), Article~98.}

The data complexity of evaluating a Boolean conjunctive query (CQ) is a
central question in database theory. The \emph{submodular width} of a query
characterises the running time of optimal combinatorial algorithms; the AGM
bound, fractional hypertree width, and the PANDA framework live in this
combinatorial setting. In parallel, \emph{fast matrix multiplication} (FMM),
which multiplies two $n \times n$ matrices in time $O(n^\omega)$ for
$\omega < 2.4$, has been used in scattered, ad-hoc results to beat the
combinatorial bounds for specific queries such as $k$-clique counting. A
unifying picture of when and how FMM helps was missing.

Abo~Khamis, Hu and Suciu close this gap. They introduce the
\emph{$\omega$-submodular width}, a refinement of submodular width
parameterised by the matrix-multiplication exponent~$\omega$. The new
parameter is never larger than the classical submodular width and equals it
when $\omega = 3$ (the FMM-free regime). Their main result, Theorem~7.1,
states that for every Boolean CQ~$Q$ and every rational $\omega$, there is
an algorithm that evaluates $Q$ in time
$\widetilde{O}\!\left(N^{\omega\text{-}\mathrm{subw}(Q)}\right)$, where $N$
is the input size and $\widetilde{O}$ hides polylogarithmic factors. This is
the first systematic statement that FMM-aided algorithms fit into the same
information-theoretic framework that yielded worst-case-optimal joins.

The proof extends the PANDA machinery. Degree configurations of the input
relations are modelled as edge-dominated polymatroids; Shannon inequalities
(monotonicity, submodularity) bound the sizes of intermediate results; and
$\omega$-submodular width is read off from a polymatroid LP. The new
ingredients are a formal calculus of \emph{matrix-multiplication
expressions} (Definition~4.2) and a corresponding \emph{variable
elimination via matrix multiplication} (Definition~4.5), which permits FMM
steps as legal moves inside an otherwise standard query plan. A generalised
Reset Lemma (Lemma~A.7), strengthening the PANDA Reset Lemma to handle
$\omega$-Shannon inequalities with proper dominance, is the technical heart
of the upper bound.

The paper does more than unify. Table~1 lists concrete query classes ---
$k$-cliques, $k$-cycles, $k$-pyramids --- for which $\omega$-submodular
width is strictly smaller than submodular width, yielding new
fastest-known algorithms. The framework also extends to disjunctive Datalog
rules, suggesting it should be useful well beyond Boolean CQs. Two
takeaways stand out: first, fast matrix multiplication --- often viewed as
an exotic primitive --- fits neatly into the polymatroid and
Shannon-inequality picture once degree configurations are made central;
second, the unification immediately yields new algorithms, not just
retrospective explanations of old ones.

\subsection*{Kapralov, Makarov, Sohler.\\
\emph{On the Adversarial Robustness of Locality-Sensitive Hashing in
Hamming Space.}\\
PACMMOD Vol.~3, No.~2 (PODS~'25), Article~102.}

Locality-sensitive hashing (LSH) is a workhorse data structure for
approximate nearest-neighbour search in high-dimensional spaces; its
variants underlie image and music retrieval, anomaly detection, clustering
and many database join primitives. The classical analysis guarantees that,
for a query chosen \emph{independently} of the data structure's random
hash functions, the probability of a false negative --- missing a point
within distance $r$ when one exists --- is at most a small parameter
$\delta$. The guarantee is fragile, however, when the same query stream
can react to earlier responses. How fast can an adaptive adversary actually
break LSH?

For LSH in Hamming space the answer turns out to be: very fast.
Kapralov, Makarov and Sohler prove (Theorem~1.1) that an adversary with
only black-box access to the LSH query interface --- no knowledge of the
hash families or seeds --- can drive the structure into a false negative
on a chosen target point using $O(\log(cr)\cdot\log(1/\delta))$ adaptive
queries, where $r$ is the LSH distance parameter and $c$ the approximation
factor. The contrast with the non-adaptive setting is dramatic: blind
random sampling needs $\Omega(1/\delta)$ trials, an exponential gap.

The attack (Algorithm~2) is an \emph{adaptive walk} starting from an
isolated target point $z$ in the dataset. The adversary samples a query at
distance $r-t$ from $z$, observes which neighbours LSH returns, and then
flips individual coordinates of the query to drift it away from $z$ while
keeping collisions inside the hash buckets that LSH actually inspects.
Lemmas~4.1, 4.5 and 4.7 control the number of hash functions involved, the
expected bucket sizes and the collision probabilities across hash
families. A Chernoff-bound analysis ties these together to show that the
walk succeeds with high probability after the stated number of steps.

The result fills a long-standing gap in the adversarial-robustness
literature. While the streaming community has produced a sequence of
attacks and defences for randomised sketches over recent years, an
analogous attack on LSH itself --- arguably the more widely deployed of the
two primitives --- had been missing. The paper makes the practical
implication explicit: any system that exposes LSH outputs to inputs that
later queries can be derived from (online retrieval, recommendation,
anomaly-detection feedback loops) cannot lean on the textbook $\delta$
guarantee. It closes by discussing robustness-aware variants and
verification as natural follow-up directions, framing the attack as an
opening shot rather than a final word.

\subsection*{Qin Zhang.\\
\emph{Robust Statistical Analysis on Streaming Data with Near-Duplicates
in General Metric Spaces.}\\
PACMMOD Vol.~3, No.~2 (PODS~'25), Article~111.}

Real streaming data --- sensor traces, image feeds, repeatedly measured
quantum states, free-text records --- is rarely clean: near-duplicates are
the norm. Two items that should be treated as the same entity are typically
only \emph{close} under some metric rather than literally equal, and an
upfront cleaning step is usually infeasible at stream rates. The
\emph{robust} versions of the two foundational streaming primitives ---
$F_0$ (the count of distinct items) and $\ell_0$-sampling (uniform sampling
from the distinct set) --- collapse near-duplicates into equivalence
classes and ask for correct behaviour on the resulting universe. Prior work
handled these robust variants only in $O(1)$-dimensional Euclidean space.

Qin Zhang lifts the picture to general metric spaces, in a single pass and
sub-linear space. The paper splits the problem into two regimes. For
\emph{well-shaped} inputs, in which equivalence classes are well separated
under the metric, two upper bounds are given: a streaming algorithm
computing a $(1\pm\varepsilon)$ approximation of robust $F_0$ in
$O\!\left(\min\{(1/\varepsilon)\sqrt{m\log(1/\varepsilon)},\,F_0\}\right)$
words (Section~2), and a robust $\ell_0$-sampler in $O(\sqrt{m})$ words
(Section~3). For arbitrary inputs, the right notion is robust $F_0$ up to
an $F_0$-ambiguity parameter $\tau$ that captures items lying in the
borderlands between equivalence classes; Section~4 gives a
$(1\pm\varepsilon)$ approximation in $O((1/\varepsilon)\sqrt{m\log m})$
words for that setting. A matching lower bound in Section~5 shows that
even a $1.1$-factor approximation requires $\Omega(\sqrt{m})$ bits,
sharply separating the general-metric setting from the classical case
where $O(\log n)$ space suffices.

The algorithmic ingredients are familiar in shape but carefully adapted to
non-Euclidean geometry: a grid partition takes the place of the cell
partition that worked in low-dimensional Euclidean space; a global
``zero-function'' threshold makes sampling consistent across the stream;
robust counts are read off as minimum-cardinality valid partitions; and
Chernoff bounds discipline deviations across the hash families used to
identify near-duplicates. It is the grid step that turns logarithmic space
into $\sqrt{m}$ space: in a general metric the number of buckets the
algorithm has to track grows polynomially in $m$ rather than only
logarithmically.

The interest of the paper is twofold. Practically, it provides the first
principled streaming primitive for robust $F_0$ and $\ell_0$-sampling in
settings whose geometry is not low-dimensional Euclidean: textual edit
distance, tree edit distance, image-feature distance, and trace distance
between quantum states. That last setting is the paper's most striking
motivation --- learning a $d$-dimensional quantum state requires
$\Omega(d/\varepsilon^2)$ noisy copies, producing a stream of
approximately-but-never-exactly-equal observations that classical
distinct-element machinery cannot consume directly. Theoretically, the
$\Omega(\sqrt{m})$ lower bound is the headline: moving from low-dimensional
Euclidean to general metrics is not a small quantitative slowdown but a
qualitative jump in how much memory robust streaming inherently needs.

\section*{References (proceedings)}

\begin{itemize}[leftmargin=*,itemsep=2pt]
  \item M.~Abo~Khamis, X.~Hu, D.~Suciu. \emph{Fast Matrix Multiplication
        meets the Submodular Width.} PACMMOD Vol.~3, No.~2 (PODS~'25),
        Article~98.
  \item M.~Kapralov, M.~Makarov, C.~Sohler. \emph{On the Adversarial
        Robustness of Locality-Sensitive Hashing in Hamming Space.}
        PACMMOD Vol.~3, No.~2 (PODS~'25), Article~102.
  \item Q.~Zhang. \emph{Robust Statistical Analysis on Streaming Data
        with Near-Duplicates in General Metric Spaces.} PACMMOD Vol.~3,
        No.~2 (PODS~'25), Article~111.
\end{itemize}

Full proceedings: \url{https://dl.acm.org/toc/pacmmod/2025/3/N2}.

\end{document}
